\subsection{Conteo por Simetría de Binomios (subsecuencias con suma alternada 0)}

\graybold{Preguntas guía:}

\begin{itemize}
\item ¿Piden contar subconjuntos o subsecuencias que cumplen una condición de suma o balance?
\item ¿Hay valores repetidos en la entrada, de forma que agruparlos por valor simplifique el conteo?
\item ¿Se puede usar C(c, par) = C(c, impar) = \pow{2}{c-1} para separar el conteo por grupo y combinar con una potencia de 2?
\end{itemize}

\graybold{Utilidad:} Contar eficientemente subconjuntos o subsecuencias que satisfacen una condición combinatoria, evitando enumerar los \pow{2}{n} subconjuntos, aprovechando la simetría de los coeficientes binomiales para reducir el problema a una subsecuencia mucho más pequeña: solo los valores distintos.

\graybold{Complejidad:} \bigO{n} por caso (la entrada ya viene no-decreciente).

\graybold{Fundamento teórico:} Dentro de un grupo de c copias iguales, elegir una cantidad par de ellas no cambia la suma alternada acumulada (aporta 0) y elegir una cantidad impar sí la cambia (aporta ±valor). Como C(c,par) = C(c,impar) = \pow{2}{c-1}, cada grupo aporta el mismo peso sin importar el modo elegido, así que el problema se reduce a contar subconjuntos de suma alternada 0 sobre los valores DISTINTOS, multiplicando el resultado por \pow{2}{n-m} al final (m = cantidad de valores distintos). En esa secuencia reducida, agrupando términos de dos en dos se demuestra que la única forma de obtener alternada 0 con más de un elemento es usando exactamente 3 valores: −1 y un par de enteros positivos consecutivos.

\graybold{Ejercicio aplicado}

(Codeforces 2246C — 0mar and Alternating Sums) Dado un arreglo no decreciente de n valores (cada uno −1 o un entero positivo), contar cuántas subsecuencias tienen suma alternada (primer término positivo, segundo negativo, tercero positivo, ...) igual a 0, módulo $\ten{9}+7$.

\graybold{I/O:} t ≤ \ten{4}, Σn ≤ 2·\ten{5} → volumen moderado; se usa readline por línea (patrón 1), adecuado porque el costo dominante está en el procesamiento \bigO{n}, no en la lectura.

\graybold{Código solución}

\begin{lstlisting}[language=Python]
import sys
input = sys.stdin.buffer.readline
 
MOD = 10**9 + 7
 
t = int(input())
out = []
for _ in range(t):
    n = int(input())
    a = list(map(int, input().split()))
 
    # el arreglo ya viene no-decreciente: colapsar a valores distintos.
    # cada grupo de c repetidos aporta un factor 2^(c-1) independiente
    # del resto (C(c,par) == C(c,impar) == 2^(c-1))
    distinct = []
    for x in a:
        if not distinct or distinct[-1] != x:
            distinct.append(x)
 
    m = len(distinct)
    has_neg1 = 1 if distinct[0] == -1 else 0   # -1 es el unico valor <= 0 posible
 
    pos = [v for v in distinct if v > 0]
    count_consec = 0
    for i in range(len(pos) - 1):
        if pos[i + 1] - pos[i] == 1:            # pares de enteros positivos consecutivos
            count_consec += 1
 
    # en la secuencia reducida, la suma alternada solo puede dar 0 con:
    # el subconjunto vacio, o exactamente {-1, v, v+1} para cada pareja
    # consecutiva (si hay un -1 disponible)
    reduced = (1 + (count_consec if has_neg1 else 0)) % MOD
    ans = reduced * pow(2, n - m, MOD) % MOD    # reincorpora los repetidos
    out.append(str(ans))
 
sys.stdout.write('\n'.join(out) + '\n')
\end{lstlisting}
